%% The following is a directive for TeXShop to indicate the main file
%%!TEX root = diss.tex

%================================
\chapter{Conclusion}
\label{ch:Conclusion}
%================================

In this thesis, we looked at the problem of generalizing methods for handling lattice walk enumeration in many dimensions and with more general constraints, as well as models with interactions. In the earlier chapters, we used refined notation and a new perspective of applying operators on functions and functional equations. We were able to derive the functional equations for \classic \dD models with ease. This allowed us to generalize the methods to derive a functional equation for a much more general class of models. This is one more step towards finding out if we can automate the solving and classification of more complex lattice walk models. However, one of the limitations of using operators is that it may be difficult to find finite equivalent operators for some rather tame-looking regions, like a wedge. It is not clear if this is due to a shortcoming of the operator notation or if this is due to some more fundamental differences between quarter plane models and those on wedges.

We were also able to extend the algebraic and obstinate kernel methods to lattice walk models with interactions, although there is still a fair amount of work that needs to be done. First, the algebraic kernel method does not work in a lot of cases of finite groups when the model has interactions, and in those cases that it does work, it does not succeed in eliminating the $\Qf{0,0}$ term. The resulting formulae still need to be worked through. At best, we hope that in those cases, if the function $\Qf{0,0}$ is D-finite, then the full generating function $Q$ is D-finite too. 

In the extension of the obstinate kernel method, we were able to find relations between $\Qf{0,0}$ and its various partially-marked versions for two nontrivial models --- the Cartesian and Diagonal models. However, we were not able to find such a relation for the Full model. It also remains to be seen if the cases with fewer transformations that preserve the admissibility of the functions $Q_i^0$ can also be tackled. The problem of solving functional equations in a closed form for more generalized models is wide open.

%================================
\section{Future Research}
%================================

In closing, the research for this thesis produced many ideas for future work, and since there were too many to pursue, we list them here for the interested reader.
\begin{itemize}
\item Algebraic kernel method.
\begin{itemize}
	\item For finite group cases, can say that $Q$ is D-finite if $\Qf{0,0}$ happens to be D-finite?
	\item Determine if the method can be extended to \highD models with interactions and/or models with larger steps.
	\item Consider annihilation of terms in the functional equations by other operators, like differential ones, and if this can be used to find out more about the generating functions.
\end{itemize}

\item Obstinate kernel method.
\begin{itemize}
	\item Can identities like $B(F\Qf{0,0}) = 0$ be explained combinatorially? The \oneD model relations would be a good place to start.
	\item Find $Q$ for the Diagonal model in terms of partially-marked versions of $Q$ and $\Qf{0,0}$.
	\item Find out if there is a similar relationship for $\Qf{0,0}$ in the Full model.
	\item Solve some of the enumeration problems with other groups, like those of order six.
\end{itemize}

\item Kernel methods in general.
\begin{itemize}
	\item Determine if the algebraic and obstinate kernel methods can be combined into one holistic approach. This can involve classifying methods on the number of substitutions made.
	\item Can the methods be extended to more generalized models?
\end{itemize}

\item Connections with other theory.
\begin{itemize}
	\item Using the Cauchy integral formula, we can rewrite extractions as complex integrals. This suggests that some of the equations can be turned into multivariable Fredholm integral equations. Perhaps methods from that area can be adapted to this one.
	\item Can we adapt Sturm-Liouville theory to lattice walk enumeration? In particular, can we find a suitable inner product for Laurent series to test operators for self-adjointness and can we expand $Q$ in terms of eigenfunctions of the lattice walk operator?
	\item Can we find asymptotics for lattice walks directly from the operator equations?
\end{itemize}
\end{itemize}
\endinput
